\chapter{第九届大唐杯B组省赛}
\fancyhead[L]{\color{H6}\kaishu\faIcon{atom}\;第九届大唐杯B组省赛}
\fancyhead[R]{\color{H6}\kaishu\rightmark\,}

\date{2022年4月10日}{}{\href{https://github.com/xubenshan/dtcup}{\textbf{才疏学浅的小熊}}}
{}
{https://www.xiaohongshu.com/user/profile/6192219d0000000021028647}{小红书}
{https://space.bilibili.com/692546609}{B站账号}



\section{单选题（共30小题，共90分）}



\begin{choice}{A}[]
    C-V2X 中，下面有关 PC5 接口表述错误的是
    \begin{tasks}(1)
        \task 是 UE 与基站之间的参考点
        \task   UE 之间的参考点
        \task  包括基于LTE的PC5
        \task 包括基于 NR 的 PC5
    \end{tasks}
\end{choice}


\begin{choice}{D}[]
    5G 基站传输规划表中，基于 IPV4“30 位掩码”对应的子网掩码应是
    \begin{tasks}(2)
        \task 255.255.255.248
        \task 255.255.255.0
        \task 255.255.255.240
        \task 255.255.255.252
    \end{tasks}
\end{choice}




\begin{choice}{C}[]
    C−V2X 中，下面有关 MEC 部署位置不正确的是
    \begin{tasks}(4)
        \task 与基站共址
        \task 骨干汇聚机房
        \task 与 RRU 共址
        \task 接入汇聚机房
    \end{tasks}
\end{choice}




\begin{choice}{B}[]
    5G NR 帧结构，常规 CP 子载波间隔不可以是
    \begin{tasks}(4)
        \task 30
        \task 20
        \task 15
        \task 240
    \end{tasks}
\end{choice}


\begin{choice}{B}[]
    5G NR 系统中，RRC INACTVE 状态，如果要恢复业务，则会触发哪条信令
    \begin{tasks}(2)
        \task RRC\;connection\; reject
        \task RRC\; resume \;request
        \task RRC \;connection\; request
        \task  RRC \;Connection \;reconfiguration
    \end{tasks}
\end{choice}

\begin{choice}{D}[]
    C-V2X 中，BSM 消息是指的是
    \begin{tasks}(2)
        \task 地图消息
        \task 路测分发交通参与者实时信息
        \task 信号灯消息
        \task  车辆基本安全类消息
    \end{tasks}
\end{choice}




\begin{choice}{C}[]
    5G NR 中，有关异系统测量的事件是
    \begin{tasks}(4)
        \task A4
        \task A6
        \task B2
        \task A1
    \end{tasks}
\end{choice}


\begin{choice}{D}[]
    在 5G 通信系统(CU DU)分离方案中，下列哪个 option 对时延的要求最为严格
    \begin{tasks}(4)
        \task  option5
        \task option1
        \task option2
        \task  option8
    \end{tasks}
\end{choice}

\begin{choice}{C}[]
    为了解决NR网络深度覆盖问题，以下哪些措施是不可取的
    \begin{tasks}(2)
        \task 采用低频段组网
        \task 增加 AAU 发射功率
        \task 增加 NR 系统带宽
        \task 使用小站提供室分覆盖
    \end{tasks}
\end{choice}


\begin{choice}{C}[]
    根据移动信道的传播特性，无线信号在自由空间的衰落情况是：传播距离每增大一倍，信号衰落增加
    \begin{tasks}(4)
        \task 9dB
        \task 3dB
        \task 6dB
        \task 2dB
    \end{tasks}
\end{choice}

\begin{choice}{D}[]
    目前工信部已经发放车联网直连通信频率，在中国车联网直连通信使用频率为
    \begin{tasks}(4)
        \task 5770-5850MHz
        \task 5855-5925MHz
        \task 5850-5925MHz
        \task 5905-5925MHz
    \end{tasks}
\end{choice}


\begin{choice}{C}[]
    5G SA 场景下，Uu 口用户面协议从上到下的次序依次是
    \begin{tasks}(2)
        \task SDAP-PHY-MAC-RLC-PDCP
        \task PDCP-PHY-MAC-RLC-SDAP
        \task SDAP-PDCP-RLC-MAC-PHY
        \task  PDCP-RLC-MAC-PHY-SDAP
    \end{tasks}
\end{choice}


\begin{choice}{A}[]
    大唐 5G 基站产品 EMB6216，最大支持安插几块 HBPOF 基带板
    \begin{tasks}(4)
        \task 6
        \task 5
        \task 3
        \task 4
    \end{tasks}
\end{choice}


\begin{choice}{A}[]
    5G NR 中哪类 SRB 不仅承载 NAS 消息还可以承载 RRC 消息，映射到 DCCH 信道
    \begin{tasks}(4)
        \task SRB1
        \task  SRB2
        \task SRB0
        \task SRB3
    \end{tasks}
\end{choice}

\begin{choice}{A}[]
    5G NR 帧结构，1 个子帧的时隙数不可能有
    \begin{tasks}(4)
        \task 3
        \task 1
        \task 4
        \task 2
    \end{tasks}
\end{choice}

\begin{choice}{D}[]
    5G NR 系统 R15 版本中，PDSCH 支持最大调制方式为
    \begin{tasks}(2)
        \task 128QAM
        \task 64QAM
        \task 512QAM
        \task 256QAM
    \end{tasks}
\end{choice}



\begin{choice}{D}[]
    5G 系统进行站内 RRC 连接重建立时，如果基站拒绝 UE 重建和新建，此时基站会发送哪条信令消息
    \begin{tasks}(4)
        \task  RRC Setup
        \task RRC establishment
        \task 不回任何消息
        \task RRC Reject
    \end{tasks}
\end{choice}

\begin{choice}{B}[]
    C−V2X 中，部署 MEC 的主要目的下面不正确的是
    \begin{tasks}(2)
        \task 为 UE 提供更快接入
        \task 传输网的负载也随之降低
        \task 降低端到端处理时延
        \task 提升用户感知
    \end{tasks}
\end{choice}


\begin{choice}{C}[]
    针对弱势交通参与者碰撞预警应用，其涉及的通信方式为
    \begin{tasks}(2)
        \task V2P/V2V
        \task V2V/V2N
        \task V2P/V2I
        \task V2P/V2N
    \end{tasks}
\end{choice}



\begin{choice}{C}[]
    在超密集组网的场景下，什么是解决边界效应的关键
    \begin{tasks}(2)
        \task 回传技术
        \task 全频谱接入
        \task 小区虚拟化
        \task 软扇区
    \end{tasks}
\end{choice}



\begin{choice}{B}[]
    5G NR 接入网构架中，CU 与 DU 划分方案中，以下选项中（\quad）要求的传输带宽最高
    \begin{tasks}(2)
        \task option2
        \task option8
        \task option4
        \task option3
    \end{tasks}
\end{choice}




\begin{choice}{C}[]
    扩展 CP 下，一个时隙的符号数是
    \begin{tasks}(4)
        \task 28
        \task 14
        \task 12
        \task 20
    \end{tasks}
\end{choice}


\begin{choice}{C}[]
    5G 方位角规划原则不正确的是
    \begin{tasks}(1)
        \task 同站两个扇区夹角不要小于 90 度
        \task 天线主瓣不要朝向湖面、江面等区域
        \task 为了覆盖更好，天线主瓣和道路的夹角越小越好
        \task 天线主瓣朝向用户密集区域
    \end{tasks}
\end{choice}


\begin{choice}{B}[]
    在 5G NR 网络中终端在初始接入过程中，基站向核心网发的第一条信息是
    \begin{tasks}(2)
        \task Uplink NAS Transfer
        \task Initial UE Message
        \task RRC Setup Complete
        \task  UE Capability Information
    \end{tasks}
\end{choice}


\begin{choice}{A}[]
    在 5G 无线网络规划中，链路预算需要与什么结合计算得到小区半径
    \begin{tasks}(2)
        \task 传播模型
        \task 系统仿真
        \task 容量规划
        \task 覆盖规划
    \end{tasks}
\end{choice}



\begin{choice}{C}[]
    5G 上行物理信号中，探测参考信号指的是
    \begin{tasks}(4)
        \task CSI-RS
        \task PTRS
        \task SRS
        \task DMRS
    \end{tasks}
\end{choice}


\begin{choice}{D}[]
    5G SSB 构成中，下面不属于 SSB 块的
    \begin{tasks}(4)
        \task SSS
        \task PBCH
        \task PSS
        \task CRS
    \end{tasks}
\end{choice}


\begin{choice}{D}[]
    C-V2X 中，下面不属于 V2X 范畴的是
    \begin{tasks}(4)
        \task V2V
        \task V2I
        \task V2N
        \task  V2C
    \end{tasks}
\end{choice}


\begin{choice}{A}[]
    C-V2X 中，基于 LTE 实现的 PC5-U 协议栈，没有下面哪个协议层
    \begin{tasks}(4)
        \task SDAP
        \task RLC
        \task PDCP
        \task MAC
    \end{tasks}
\end{choice}


\begin{choice}{B}[]
    eMBB 场景下，不属于影响 5G 系统端到端时延的维度是
    \begin{tasks}(2)
        \task 传输网(Transport)延时
        \task ftp 服务器延时
        \task 无线接入网(RAN)延时
        \task 核心网(Core)延时
    \end{tasks}
\end{choice}



\section{多选题（共20小题，共80分）}

\begin{choice}{\;ABC\;}[]
    关于 5G 中 QoS flow 与 DRB 关系下面说法正确的是
    \begin{tasks}(2)
        \task 多个 QoS flow 映射一个 DRB
        \task DRB 保证 QoS flow 空口质量
        \task 一个 QoS flow 映射一个 DRB
        \task 一个 QoS flow 可以映射多个 DRB
    \end{tasks}
\end{choice}



\begin{choice}{\;ABCD\;}[]
    5G RRC 连接重建触发有如下原因
    \begin{tasks}(4)
        \task 切换失败
        \task 重配失败
        \task 完保失败
        \task 基站检测 RLF
    \end{tasks}
\end{choice}


\begin{choice}{\;BC\;}[]
    根据不同应用对通信频率和时延的不同需求，将 17 个一期应用分为两大类，分别为
    \begin{tasks}(2)
        \task 低时延、低频率
        \task 高时延、低频率
        \task 低时延、高频率
        \task 高时延、高频率
    \end{tasks}
\end{choice}

\begin{choice}{\;ABCD\;}[]
    5G 通信系统中，哪些关键技术起到了缩短时延的作用
    \begin{tasks}(2)
        \task RRC-inactive 状态
        \task MEC 网络架构
        \task 自包含帧结构
        \task Shorter-TTI
    \end{tasks}
\end{choice}

\begin{choice}{\;ACD\;}[]
    下面那些接口属于 C−2X 的 PC5 接口
    \begin{tasks}(4)
        \task 车与 RSU（UE 级别）
        \task 车与基站
        \task 车与车之间
        \task 车与人
    \end{tasks}
\end{choice}

\begin{choice}{\;ABCD\;}[]
    5G EN-DC 下可以建立哪些 SRB 类型
    \begin{tasks}(4)
        \task  SRB2
        \task SRB1
        \task SRB3
        \task SRB0
    \end{tasks}
\end{choice}


\begin{choice}{\;BCD\;}[]
    5G 在移动通信系统中关于阴影衰落的描述，正确的是
    \begin{tasks}(1)
        \task 阴影衰落又叫快衰落
        \task 电波传播路径上遇有高大建筑物、树林、地形起伏等障碍物的阻挡，就会产生电磁场的阴影衰落
        \task 阴影衰落又叫慢衰落
        \task 这种慢衰落的规律，其中值变动服从对数正态分布
    \end{tasks}
\end{choice}

\begin{choice}{\;BCD\;}[]
    5G NR 中，RRC 状态包括哪些
    \begin{tasks}(2)
        \task RM REGISTERED
        \task RRC IDLE
        \task RRC INACTIVE
        \task RRC CONNECTED
    \end{tasks}
\end{choice}

\begin{choice}{\;AD\;}[]
    下面哪些属于 5G 上行信号
    \begin{tasks}(4)
        \task SRS
        \task CSI-RS
        \task  SSS
        \task PUSCH\;DMRS
    \end{tasks}
\end{choice}


\begin{choice}{\;BCD\;}[]
    有关 GPS 天线安装，下面要求正确的是
    \begin{tasks}(1)
        \task 两个 GPS 天线间距大于 1 米
        \task 天线竖直向上的视角应大于 120 度
        \task 在塔上安装 GPS 天线，应使 GPS 天线位于铁塔最面向南方一角并保持与铁塔 2 米间距
        \task 周围对天线的遮挡不超 30 度
    \end{tasks}
\end{choice}


\begin{choice}{\;BCD\;}[]
    根据合作式智能运输系统车用通信系统应用层及应用数据交互标准，选取 17 个典型应用作为一期应用，这些应用可分为不同类别，分别为
    \begin{tasks}(4)
        \task 自动驾驶
        \task 信息服务
        \task 安全
        \task  效率
    \end{tasks}
\end{choice}


\begin{choice}{\;ABCD\;}[]
    以下关于 5G 网络架构的发展方向，说法正确的是
    \begin{tasks}(1)
        \task 支持面向客户的业务模式
        \task 同时支持各种差异化场景
        \task 支持业务的快速建立和修改
        \task 虚拟化、组件化、可编排是 5G 网络架构特点
    \end{tasks}
\end{choice}


\begin{choice}{\;ACD\;}[]
    大唐 5G 设备安装时，以下 BBU 竖装规范要求，描述正确的是
    \begin{tasks}(1)
        \task BBU竖装时，前方需留有 800mm 的维护通道
        \task  BBU 竖装时，机柜后部不支持靠墙安装，与墙的可距离≤800mm
        \task BBU 竖装支持并柜或并架安装
        \task BBU 竖装双排或多排机柜安装时，机柜前后间隔≥800mm，同时需要同风向安装；
    \end{tasks}
\end{choice}

\begin{choice}{\;BCD\;}[]
    NR 系统中，关于路测弱覆盖原因分析，以下描述正确的是哪些项
    \begin{tasks}(1)
        \task 传输资源拥塞
        \task 设备问题：设备工作状态异常，AAU 硬件故障等
        \task 可能是环境原因，如建筑物遮挡
        \task 切换参数配置，导致无法及时切换
    \end{tasks}
\end{choice}

\begin{choice}{\;ACD\;}[]
    关于越区覆盖，以下描述正确的是哪些项
    \begin{tasks}(1)
        \task 越区覆盖可能形成“孤岛”，导致掉话
        \task 对于 NR 系统而言越区覆盖不会对其邻近小区形成干扰
        \task 越区覆盖可能导致越区覆盖基站与周边基站负荷不均。周边基站资源利用率低
        \task 越区覆盖可能导致 UE 发射功率受限，引起接入失败
    \end{tasks}
\end{choice}

\begin{choice}{\;ABD\;}[]
    现阶段，非独立组网架构(NSA)包括以下哪几种 option(选项)
    \begin{tasks}(4)
        \task  Option4/4a
        \task Option3/3a/3x 
        \task Option2
        \task Option7/7a/7x
    \end{tasks}
\end{choice}

\begin{choice}{\;ABCD\;}[]
    大唐 5G 基站开通时，板卡不可用的原因可能为
    \begin{tasks}(2)
        \task 时钟问题
        \task 板卡规划错误
        \task  插入板卡类型不对 
        \task 板卡上处理器加载失败
    \end{tasks}
\end{choice}

\begin{choice}{\;BD\;}[]
    5G NR 中，3GPP 主要指定了两个频率范围，分别是
    \begin{tasks}(4)
        \task  FR4 
        \task FR2 
        \task FR3 
        \task  FR1
    \end{tasks}
\end{choice}

\begin{choice}{\;ABCD\;}[]
    UDN 面临的主要问题是
    \begin{tasks}(2)
        \task 小区间干扰问题
        \task 站址难选择问题
        \task 回传部署问题
        \task 频繁切换重选问题
    \end{tasks}
\end{choice}

\begin{choice}{\;ABD\;}[]
    5G 高频段的通信特点是
    \begin{tasks}(2)
        \task 距离短
        \task 穿透弱
        \task 天线阵尺寸大
        \task 损耗大
    \end{tasks}
\end{choice}



\section{判断题（共10小题，共30分）}

\begin{choice}{\;正确\;}[]
    5G 系统，调制编码等级越高，单个 RE 承载的 bit 数越多。
\end{choice}


\begin{choice}{\;错误\;}[]
    5G 测试中，如果发现 UE 连接状态时弱覆盖，该问题肯定与邻区漏配无关。
\end{choice}

\begin{choice}{\;错误\;}[]
    C-V2X 只能采用 NR 实现，不能采用 LTE 技术实现。
\end{choice}

\begin{choice}{\;正确\;}[]
    5G NR 小区全局标识(NCGI)的组成部分包括 MCC\;MNC\;gNB ID 和 Cell ID。
\end{choice}

\begin{choice}{\;正确\;}[]
5G 系统中，PCI mod3/4 相同对 SS-SINR 影响很小，基本不用考虑该因素。
\end{choice}

\begin{choice}{\;正确\;}[]
    大唐 5G 网络运维过程中，BBU 侧发送光功率最小值不能低于-10dBm。

\end{choice}


\begin{choice}{\;错误\;}[]
    5G RRC 连接建立是建立 SRB2。

\end{choice}

\begin{choice}{\;正确\;}[]
    C-V2X 网络架构中的管即管道，负责终端采集数据的回传和控制信息的分发。

\end{choice}

\begin{choice}{\;正确\;}[]

    5G SRB0 映射的逻辑信道是 CCCH。
\end{choice}


\begin{choice}{\;正确\;}[]
    若两台主机的 IP 地址分别与它们的子网掩码相“与”后的结果相同，则说明这两台主机在同一子网中。

\end{choice}


